% chapters/implementation/subsystems.tex
% Sección 5.2: Desarrollo por subsistemas (~4-5 páginas)
% ============================================================================
%
% GUÍA: Describir la implementación de cada subsistema funcional.
%       No repetir el diseño del cap. 4; centrarse en CÓMO se implementó:
%       decisiones técnicas, problemas encontrados, soluciones aplicadas.
%       Incluir fragmentos de código representativos si aportan valor.
%
%   Para cada subsistema, usar \subsection (sin numerar en el índice):
%
%   \subsection{Modelos de datos y migraciones}
%   - Implementación de los modelos Django.
%   - Estrategia de migraciones.
%   - Soft delete y anonimización (RGPD).
%   - Ejemplo de código relevante (si procede):
%     \PythonCode[cod:model_example]{Modelo de ejemplo}
%                {Ejemplo del modelo ExamInstance con campo de estado}
%                {codes/exam_instance_model.py}{1}{25}{1}
%   TODO: Copiar fragmentos de código representativos al directorio codes/.
%
%   \subsection{Sistema de autenticación}
%   - Implementación del patrón Strategy.
%   - Cómo se registra y carga el plugin activo.
%   - Detalles de JWT (librerías usadas, configuración de expiración).
%   - Detalles de sesión (configuración de cookies seguras).
%
%   \subsection{Sistema de permisos}
%   - Implementación del modelo de bloques funcionales.
%   - Evaluación de la cadena de permisos.
%   - Integración con DRF (permission classes personalizadas).
%
%   \subsection{Gestión académica y CRUD}
%   - ViewSets y routers de DRF.
%   - Serializers con validaciones.
%   - Importación masiva CSV/JSON.
%   - Filtrado, ordenación y paginación.
%
%   \subsection{Pipeline de ingesta y OCR}
%   - Implementación del SFTP watcher (inotify, polling, etc.).
%   - Tareas Celery para procesamiento asíncrono.
%   - Integración con MinIO para almacenamiento.
%   - Implementación del plugin OCR.
%   - Algoritmo de identificación por similitud.
%   - Gestión de errores y reintentos.
%
%   \subsection{Corrección y anotaciones}
%   - Asignación automática de correctores (motor de reglas).
%   - Almacenamiento de anotaciones multimedia en MinIO.
%   - Control de concurrencia optimista (version field).
%   - Cálculo de calificaciones.
%
%   \subsection{Ciclo de vida del examen}
%   - Implementación de la máquina de estados.
%   - Transiciones validadas y acciones asociadas.
%   - Publicación de notas y ventana de revisión.
%
%   \subsection{Notificaciones y correo electrónico}
%   - Notificaciones internas (modelo + endpoint).
%   - Envío de correos con Celery (asíncrono).
%   - Plantillas multiidioma con ficheros JSON.
%
%   \subsection{Visualización y protección de documentos}
%   - Rasterización con marca de agua.
%   - Proxy de MinIO con control de acceso.
%   - Cifrado del payload.
% ============================================================================

% TODO: Redactar el desarrollo por subsistemas.
%       Incluir fragmentos de código relevantes en codes/ y referenciarlos.
